\section{Clase 16 - Pr\'actica 6 (Resoluci\'on num\'erica de ecuaciones no lineales) y Pr\'actica 7 (Interpolaci\'on)}

\subsection{Orden de convergencia}

Sea $\{a_{n}\}_{n \in \mathbb{N}}$ una sucesi\'on de n\'umeros reales, $a_{n} \rightarrow_{n \to \infty} l, l \in \mathbb{R}$. Decimos que $a_{n}$ converge a $l$ con orden al menos $p$ si existen $C > 0, n_{0} \in \mathbb{N}$ tales que $\frac{|a_{n + 1} - l|}{|a_{n} - l|} \leq C\ \forall n \geq n_{0}$.\\

Decimos que $a_{n}$ converge a $l$ con orden exactamente $p$ si converge con orden al menos $p$ y $\forall \epsilon > 0$ $\frac{|a_{n + 1} - l|}{|a_{n} - l|^{p - \epsilon}} \rightarrow_{n \to +\infty} 0$.\\

\underline{Ra\'ices m\'ultiples}\\

Sea $r$ tal que $f(r) = 0, ..., f^{(p - 1)}(r) = 0, f \in C^{p}$. Adem\'as, $f(x) = (x - r)^{p}g(x), g(r) \neq 0, g \in C^{1}$\\

\begin{itemize}
    \item Para Newton-Raphson:

    $$x_{n + 1} = x_{n} - \frac{(x_{n} - r)^{p}g(x_{n})}{p(x_{n} - r)^{p - 1}g(x_{n}) + (x_{n} - r)^{p}g'(x_{n})}$$

    $$e_{n + 1} = e_{n} - \frac{e_{n}^{p}g(x_{n})}{pe_{n}^{p - 1}g(x_{n}) + e_{n}^{p}g'(x_{n})} = e_{n} - \frac{e_{n}g(x_{n})}{pg(x_{n}) + e_{n}g'(x_{n})}$$

    $$\Rightarrow e_{n + 1} = e_{n}\frac{(p - 1)g(x_{n})}{pg(x_{n}) + e_{n}g'(x_{n})} + e_{n}^{2}\frac{g'(x_{n})}{pg(x_{n}) + e_{n}g'(x_{n})}$$

    Si $x_{n} \rightarrow r$ y $p \neq 1$,

    $$\frac{(p - 1)g(x_{n})}{pg(x_{n}) + e_{n}g'(x_{n})} \rightarrow \frac{(p - 1)g(r)}{pg(r)} = \frac{p - 1}{p} \neq 0$$

    y la convergencia es lineal.\\

    \item Para punto fijo:\\

    Si $g \in C^{1}$ y $x_{n + 1} = g(x_{n}) \rightarrow r$,

    $$e_{n + 1} = x_{n + 1} - r = g(x_{n}) - g(r) = g'(\zeta_{n})(x_{n} - r) = g'(\zeta_{n})e_{n}$$

    Si $e_{n} \neq 0, \frac{e_{n + 1}}{e_{n}} = g'(\zeta_{n})) \rightarrow g'(r)$. Si $g'(r) \neq 0$, la convergencia es lineal. Si $g \in C^{p}, g^{(k)} = 0$ para $1 \leq k \leq p, g^{(p)} \neq 0$,

    $$e_{n + 1} = x_{n + 1} - r = g(x_{n}) - g(r) = g'(r)(x_{n} - r) + \frac{g''(r)}{2}(x_{n} - r)^{2} + ... + \frac{g^{(p)}(\zeta_{n})}{p!}(x_{n} - r)^{p}$$

    $$\Rightarrow e_{n + 1} = \frac{g^{(p)}(\zeta_{n})}{p!}e_{n}^{p} \Rightarrow \frac{e_{n + 1}}{e_{n}} \rightarrow \frac{g^{(p)}(r)}{p!} \neq 0$$

    y el m\'etodo tiene orden $p$.
\end{itemize}

\subsection{Sistemas de ecuaciones no lineales}

$$
    \begin{cases}
        f_{1}(x, y) = 0\\
        f_{2}(x, y) = 0
    \end{cases}
$$\\

Llamamos $F(x, y) = (f_{1}(x, y), f_{2}(x, y)),\ F: \mathbb{R}^{2} \rightarrow \mathbb{R}^{2}$. Buscamos $\Vec{r} = (r_{1}, r_{2})$ tal que $F(\Vec{r}) = (0, 0)$.\\

\underline{Newton-Raphson}\\

Recordemos la ecuaci\'on del plano tangente en $(x_{0}, y_{0}, f(x_{0}, y_{0}))$:

$$z - f(x_{0}, y_{0}) = \frac{\partial f}{\partial x}(x_{0}, y_{0})(x - x_{0}) + \frac{\partial f}{\partial y}(x_{0}, y_{0})(y - y_{0})$$\\

Propongo $(x_{1}, y_{1})$ tal que:

$$ 0 = \frac{\partial f_{i}}{\partial x}(x_{0}, y_{0})(x_{1} - x_{0}) + \frac{\partial f_{i}}{\partial y}(x_{0}, y_{0})(y_{1} - y_{0}) + f_{i}(x_{0}, y_{0}),\ i = 1, 2$$\\

Reordeno:

$$ -\frac{\partial f_{i}}{\partial x}(x_{0}, y_{0})x_{1} - \frac{\partial f_{i}}{\partial y}(x_{0}, y_{0})y_{1} = -\frac{\partial f_{i}}{\partial x}(x_{0}, y_{0})x_{0} - \frac{\partial f_{i}}{\partial y}(x_{0}, y_{0})y_{0} + f_{i}(x_{0}, y_{0}),\ i = 1, 2$$\\

En forma matricial:

$$    -
    \begin{pmatrix}
        \frac{\partial f_{1}}{\partial x}(x_{0}, y_{0}) & \frac{\partial f_{1}}{\partial y}(x_{0}, y_{0})\\
        \frac{\partial f_{2}}{\partial x}(x_{0}, y_{0}) & \frac{\partial f_{2}}{\partial y}(x_{0}, y_{0})
    \end{pmatrix}
    \begin{pmatrix}
        x_{1}\\
        y_{1}
    \end{pmatrix}
    = -
    \begin{pmatrix}
        \frac{\partial f_{1}}{\partial x}(x_{0}, y_{0}) & \frac{\partial f_{1}}{\partial y}(x_{0}, y_{0})\\
        \frac{\partial f_{2}}{\partial x}(x_{0}, y_{0}) & \frac{\partial f_{2}}{\partial y}(x_{0}, y_{0})
    \end{pmatrix}
    \begin{pmatrix}
        x_{0}\\
        y_{0}
    \end{pmatrix}
    +
    \begin{pmatrix}
        f_{1}(x_{0}, y_{0})\\
        f_{2}(x_{0}, y_{0})
    \end{pmatrix}
$$\\

Si

$$J_{n} = Jac\ F(x_{n}, y_{n}) = \begin{pmatrix}
        \frac{\partial f_{1}}{\partial x}(x_{n}, y_{n}) & \frac{\partial f_{1}}{\partial y}(x_{n}, y_{n})\\
        \frac{\partial f_{2}}{\partial x}(x_{n}, y_{n}) & \frac{\partial f_{2}}{\partial y}(x_{n}, y_{n})
    \end{pmatrix}$$

y $J_{n}^{-1}$ su inversa, el m\'etodo es:

\begin{itemize}
    \item Tomo $x_{0} \in \mathbb{R}^{2}$ que aproxima a $\Vec{r}$

    \item Para $n = 0, 1, 2, ...$, $\Vec{x}_{n + 1} = \Vec{x}_{n} - J_{n}^{-1}F(\Vec{x}_{n})$\\

    \textit{Observaci\'on}: $J_{n}^{-1}F(\Vec{x}_{n}) = \Vec{v}_{n} \iff J_{n}\Vec{v}_{n} = F(\Vec{x}_{n})$. Es decir, $\Vec{v}_{n}$ es la soluci\'on del sistema lineal $A\Vec{x} = \Vec{b}$, con $A = J_{n}$ y $\Vec{b} = F(\Vec{x}_{n})$.\\
\end{itemize}

\underline{Punto fijo}

$$F(x, y) = (0, 0) \Rightarrow G(x, y) = (x, y)$$

El m\'etodo nos queda:

\begin{itemize}
    \item Tomo $x_{0} \in \mathbb{R}^{2}$ que aproxima a $\Vec{r}$

    \item Para $n = 0, 1, 2, ...$, $\Vec{x}_{n + 1} = G(\Vec{x}_{n})$\\
\end{itemize}

\textit{Ejemplo}: Aproximar por el m\'etodo de Newton las $x$ del sistema no lineal

$$
    \begin{cases}
        x^{2} + y^{2} = 10\\
        x^{2} - y^{2} = 1
    \end{cases}
$$\\

Tenemos $f_{1}(x, y ) = x^{2} + y^{2} - 10$ y $f_{2}(x, y) = x^{2} - y^{2} - 1$. Entonces

$$Jac\ F(x, y) = \begin{pmatrix}
        2x & 2y\\
        2x & -2y
    \end{pmatrix}$$\\

\underline{\textbf{Interpolaci\'on}}\\

Partimos de una tabla de datos

\begin{table}[h]
    \centering
    \begin{tabular}{|c|c|c|c|c|c|}
        \hline
        $x$ & $x_{0}$ & $x_{1}$ & $x_{2}$ & $...$ & $x_{n}$\\ \hline
        $y$ & $y_{0}$ & $y_{1}$ & $y_{2}$ & $...$ & $y_{n}$\\ \hline
    \end{tabular}
\end{table}

con $x_{0}, x_{1}, \cdots , x_{n}$ distintos. Buscamos un polinomio de grado $\leq n$ que pase por esos puntos.\\

\underline{Teorema}: Dados $n + 1$ valores distintos $x_{0}, x_{1}, \cdots , x_{n}$ y $n + 1$ valores cualesquiera $y_{0}, y_{1}, \cdots , y_{n}$, existe un polinomio $P_{n}$ de grado $\leq n$ que satisface $P_{n}(x_{i}) = y_{i}, \ i = 0, \cdots , n$.\\

Si $y_{i} = f(x_{j})$ para cierta funci\'on $f$, decimos que $P_{n}$ interpola a $f$ en $x_{0}, ... , x_{n}$.\\

Vamos a ver tres maneras de armar $P_{n}$: coeficientes indeterminados, forma de Lagrange y forma de Newton (tabla de diferencias finitas).\\

\subsection{Coeficientes indeterminados}

$$P_{n}(x) = a_{0} + a_{1}x + a_{2}x^{2} + ... + a_{n}x^{n}$$

y resuelvo un sistema lineal para hallar $a_{0}, ... ,a_{n}$.\\

$$P_{n}(x_{0}) = a_{0} + a_{1}x_{0} + a_{2}x_{0}^{2} + ... + a_{n}x_{0}^{n} = y_{0}$$
$$ \vdots $$
$$P_{n}(x_{n}) = a_{0} + a_{1}x_{n} + a_{2}x_{n}^{2} + ... + a_{n}x_{n}^{n} = y_{n}$$\\

En forma matricial (matriz de Vandermonde):

$$
    \begin{pmatrix}
        1 & x_{0} & x_{0}^{2} & \cdots & x_{0}^{n}\\
        \vdots & \vdots & \vdots & \cdots & \vdots\\
        \vdots & \vdots & \vdots & \cdots & \vdots\\
        1 & x_{n} & x_{n}^{2} & \cdots & x_{n}^{n}\\
    \end{pmatrix}
    \begin{pmatrix}
        a_{0}\\
        a_{1}\\
        \vdots\\
        a_{n}\\
    \end{pmatrix}
    =
    \begin{pmatrix}
        y_{0}\\
        y_{1}\\
        \vdots\\
        y_{n}\\
    \end{pmatrix}
$$\\

\textit{Ejemplo}:

\begin{table}[h]
    \centering
    \begin{tabular}{|c|c|c|c|}
        \hline
        $x$ & -1 & 0 & 1\\ \hline
        $y$ & 12 & 5 & 4\\ \hline
    \end{tabular}
\end{table}

Tenemos 3 puntos, luego $n = 2 \rightarrow P_{2}(x) = a_{0} + a_{1}x + a_{2}x^{2}$, armamos la matriz de Vandermonde y resolvemos:

$$
    \begin{pmatrix}
        1 & -1 & 2\\
        1 & 0 & 0\\
        1 & 1 & 1\\
    \end{pmatrix}
    \begin{pmatrix}
        a_{0}\\
        a_{1}\\
        a_{2}\\
    \end{pmatrix}
    =
    \begin{pmatrix}
        12\\
        5\\
        4\\
    \end{pmatrix}
$$\\

$$ \Rightarrow P_{2}(x) = 5 - 4x + 3x^{2}$$

\subsection{Forma de Lagrange}

$$P_{n}(x) = y_{0}l_{0}(x) + y_{1}l_{1}(x) + y_{2}l_{2}(x) + \cdots + y_{n}l_{n}(x)$$\\

con $l_{0}, l_{1}, \cdots, l_{n}$ polinomios tales que $l_{i}(x_{j}) = \delta_{ij}$.\\

Expl\'icitamente:

$$l_{j}(x) = \frac{(x - x_{0})(x - x_{1})(x - x_{j - 1})(x - x_{j + 1})\cdots(x - x_{n})}{(x_{j} - x_{0})(x_{j} - x_{1})(x_{j} - x_{j - 1})(x_{j} - x_{j + 1})\cdots(x_{j} - x_{n})}$$\\

con $\{l_{0}, \cdots, l_{n}\}$ base de Lagrange.\\

\textit{Ejemplo anterior}:

$$P_{2}(x) = y_{0}l_{0}(x) + y_{1}l_{1}(x) + y_{2}l_{2}(x) = 12l_{0}(x) + 5l_{1}(x) + 4l_{2}(x)$$

$$l_{0}(x) = \frac{(x - 0)(x - 1)}{(-1 + 0)(-1 - 1)} = \frac{1}{2}x(x - 1)$$

$$l_{1}(x) = \frac{(x + 1)(x - 1)}{(0 + 1)(0 - 1)} = -(x + 1)(x - 1)$$

$$l_{2}(x) = \frac{(x + 1)(x - 0)}{(1 + 1)(1 - 0} = \frac{1}{2}(x + 1)x$$

$$\Rightarrow P_{2}(x) = 12\frac{1}{2}x(x - 1) - 5(x + 1)(x - 1) + 4\frac{1}{2}(x + 1)x = 3x^{2} - 4x + 5$$

\subsection{Forma de Newton - Tabla de diferencias divididas}

\begin{table}[h]
    \centering
    \begin{tabular}{c|ccc}
        $x_{i}$ & $f[\cdot]$ & $f[\cdot, \cdot]$ & $f[\cdot, \cdot, \cdot]$\\ \hline
        $x_{0}$ & $f(x_{0})$ &  & \\
        &&&\\
        $x_{1}$ & $f(x_{1})$ & $f[x_{0}, x_{1}] = \frac{f(x_{1}) - f(x_{0})}{x_{1} - x_{0}}$ & \\
        &&&\\
        $x_{2}$ & $f(x_{2})$ & $f[x_{1}, x_{2}] = \frac{f(x_{2}) - f(x_{1})}{x_{2} - x_{1}}$ & $f[x_{0}, x_{1}, x_{2}] = \frac{f[x_{1}, x_{2}] - f[x_{0}, x_{1}]}{x_{2} - x_{0}}$\\
    \end{tabular}
\end{table}

$$P_{n}(x) = f[x_{0}] + f[x_{0}, x_{1}](x - x_{0}) + f[x_{0}, x_{1}, x_{2}](x - x_{0})(x - x_{1}) + \cdots + f[x_{0}, \cdots, x_{n}](x - x_{0})...(x - x_{n - 1})$$\\

\underline{Definici\'on}: Diferencias divididas

\begin{itemize}
    \item $f[x_{i}] = f(x_{i})$, de orden 0.
    \item $f[x_{i}, x_{i + 1}, \cdots, x_{i + m}] = \frac{f[x_{i + 1}, \cdots, x_{i + m}] - f[x_{i}, \cdots, x_{i + m - 1}]}{x_{i + m} - x_{i}}$, de orden $m$.\\
\end{itemize}

\newpage
\textit{Ejemplo anterior}:

\begin{table}[h]
    \centering
    \begin{tabular}{c|ccc}
        $x_{i} $ & $f[\cdot]$ & $f[\cdot, \cdot]$ & $f[\cdot, \cdot, \cdot]$\\ \hline
        $x_{0} = -1$ & 12 &  & \\
        &&&\\
        $x_{1} = 0$ & 5 & $\frac{5 - 12}{0 - (-1)} = -7$ & \\
        &&&\\
        $x_{2} = 1$ & 4 & $\frac{4 - 5}{1 - 0} = -1$ & $\frac{-1 - (-7)}{1 - (-1)} = 3$\\
    \end{tabular}
\end{table}

$$\Rightarrow P_{2}(x) = 12 - 7(x - (-1)) + 3(x - 0)(x - (-1)) = 3x^{2} - 4x + 5$$\\

\textit{Observaci\'on}:

\begin{itemize}
    \item Si $y_{i} = f(x_{i})$ y $f$ es un polinomio de grado $\leq n \Rightarrow P_{n} = f$.
    \item Hay infinitos polinomios de grado $\geq n + 1$ que interpolan en $n$ puntos. Por ejemplo: $P_{n + 1}(x) = P_{n}(x) + \alpha(x - x_{0})(x - x_{1})...(x - x_{n}),\ \alpha \in \mathbb{R}$\\
\end{itemize}

\underline{Error de interpolaci\'on}: $f:[a, b] \rightarrow \mathbb{R},\ x_{0}, \cdots, x_{n} \in [a, b]$ todos distintos, $x \in [a, b]$:

$$f(x) - P_{n}(x) = \frac{f^{(n + 1)}(\theta)}{(n + 1)!}\prod_{j = 0}^{n} (x - x_{j}),\ \theta \in [a, b]$$